\section{Introduction}
\label{sec:introduction}

The United States intercity bus market has been in structural decline for four decades. Following deregulation under the Bus Regulatory Reform Act of 1982 \citep{BusRegulatoryReform1982}, operators systematically withdrew from low-density corridors as competitive entry on high-density routes squeezed margins. The 2020 Greyhound bankruptcy and subsequent network restructuring eliminated service to hundreds of communities that had no alternative \citep{Schwieterman2019}. What remains is a patchwork: a handful of profitable premium routes (New York--Washington, Los Angeles--San Francisco), and a near-total absence of scheduled intercity bus service across the continental interior.

The cause is not insufficient demand. It is insufficient infrastructure. Intercity bus profitability depends on a property that current US interstate highways systematically lack: schedule reliability. When the 95th percentile travel time on a corridor is 86\% longer than the free-flow travel time --- the measured Planning Time Index (PTI) of 1.86 on I-80 between the Midwest and the coasts \citep{FHWA_freight2023} --- a bus operator faces an impossible choice. Padding the schedule by 86\% produces a service so slow it cannot compete with driving. Running an optimistic schedule produces chronic lateness that destroys ridership. The rational response is market exit, and that is precisely what has occurred.

Interstate 2.0 managed freight lanes change this calculus at the infrastructure level. T1 Primary Artery managed lanes are designed to a PTI target of $\leq$ 1.15, sustained at 65 mph, separated from general-purpose traffic \citep{ROUTE_E2}. This is not a marginal improvement; it is a phase transition. With PTI 1.15, a bus operator can advertise a departure time and guarantee 95\% on-time arrival within 15\% of the scheduled travel time. That guarantee --- the standard Amtrak and FlixBus EU use to build brand loyalty --- is currently unachievable on most US interstate corridors.

This paper, the second in Track F of the Interstate 2.0 research series, builds directly on F.1 \citep{ROUTE_F1}, which established the hub network: 9 T1/T1 diamond hubs and approximately 50 T1/T2 regional stops, serving 12.4 million transit-dependent Americans within 30 miles of a hub at a hub increment cost of \$2 billion (0.8\% of total I2.0 investment). F.1 showed that the physical infrastructure for an intercity transit layer exists wherever T1/T1 hubs are built. F.2 shows what that transit layer produces: actual travel times, operator economics, and equity outcomes on 12 specific corridors.

The central findings are as follows:

\begin{enumerate}
  \item \textbf{Travel time.} I2.0 T1 express bus achieves an average scheduled speed of 62 mph across 12 analyzed corridors, yielding travel times 28--45\% shorter than current bus alternatives on every corridor (Section~\ref{sec:travel-time}).

  \item \textbf{Modal parity.} On corridors where Amtrak service exists, T1 bus is competitive with or faster than rail; on the five corridors with no current bus or Amtrak service, T1 bus creates connectivity that does not otherwise exist (Section~\ref{sec:travel-time}).

  \item \textbf{Operator viability.} At a fare of \$0.12--0.15/mile --- comparable to Greyhound economy pricing --- T1 bus routes achieve positive operating margins on corridors with population sufficient to sustain a twice-daily round trip, without per-passenger operating subsidy (Section~\ref{sec:economics}).

  \item \textbf{Equity reach.} In-between communities along T1 corridors --- agricultural service centers, county seats, and small cities not served by air or rail --- disproportionately score high on C3 (Economic Opportunity Access deficit), and gain the most from T1 bus access to regional employment, healthcare, and educational centers (Section~\ref{sec:equity}).
\end{enumerate}

The paper proceeds as follows. Section~\ref{sec:background} reviews intercity bus economics and the infrastructure conditions for viable scheduled service. Section~\ref{sec:methods} specifies the travel time model and ridership estimation approach. Section~\ref{sec:travel-time} presents the 12-corridor travel time analysis. Section~\ref{sec:economics} develops the operator cost model and subsidy logic for marginal corridors. Section~\ref{sec:equity} quantifies the in-between community impact. Section~\ref{sec:conclusion} states policy recommendations and limitations.
